\chapter{Probability}

Probability is a measure of chance of occurance of a phenomenon.
The word ``probability`` may be used to represent ``the degree of belief`` 
of a person making a statement on the proposition. It is used in the sense 

\begin{definition}[Mutually Exclusive Events]
    Several events $E_1, E_2, \ldots, E_n$ in relation to a 
    random experiment are said to be \textbf{mutually exclusive} (or disjoint) if any 
    two of them cannot occur simultaneously, which means everytime the 
    experiment is performed is $E_i \cap E_j = \varnothing$ for all $i \neq j, i < j = 1,\ldots n$.
\end{definition}
\begin{remark}
    Mutually exclusive events are not necessarily independent, or vice versa.
\end{remark}

\begin{definition}[Exhaustive Events]
    A sequence of events $E_1, E_2, \ldots, E_n$ from a random experiment are said to be 
    exhaustive events if any of them must necessarily occur, everytime the experiment 
    is performed that is $\bigcup^n_{i=1} E_i = \Omega$.
\end{definition}

\begin{example}
    A fair coin is tossed twice, the sample space is 
    \[
        \Omega = \{ HH, TH, HT, TT \}.
    \]
    We define two events as following:
    \[
        A = \{ \text{Head appeared twice}\} = \{ HH \},
    \]
\end{example}

\section{The classical definition of probability}

If a random experiment can result in $N$ finite number of mutually exclusive, exhaustive and 
equally likely cases and $\hash A$ of them are favorable to the occurance of the event $A$, 
then the probability of occurance of $A$ is 
\begin{equation}
    \mathbb{P}(A) = \frac{\hash A }{N}.
\end{equation}
Since $0 \leq \hash A \leq N \implies 0 \leq \mathbb{P}(A) \leq 1$. By classical definition 
of probability of an event is a rational number lies between 0 and 1. But in general,
probability is a real number between 0 and 1.

However, the classical probability has its limitation:
\begin{itemize}
    \item It is assumed that all the cases are equally likely. The classical definition
    of probability is found useful when applied to the outcomes of the games of chance. If the 
    outcomes of experiment are not equally likely then this definition is not applicable.
    \item The definition of classical probability may breaks down if the number of 
    all possible outcomes is infinite.
    \item In real life it is not easy to determine the outcomes as equally likely.
\end{itemize}

Suppose $E$ is an event of a random experiment. Suppose it is possible to repeat the 
experiment in a large number of times under essentially similar conditions.

\begin{definition}[Principle of randomness]
    The limits are invariant with respect to the choice of every subsequence of the 
    collectives, which is independent of the attribute under consideration. This 
    implies that all outcomes should have either equal, or known, probability. So that 
    each individual outcomes unpredictable despite there is a known probability distribution.
\end{definition}

\section{Axiomatic Explanation}

\begin{axiom}[Kolmogorov's axioms]
Let $\Omega$ be the sample space of a random experiment and $F$ be a $\sigma$-field 
of events of $\Omega$. A set function $\mathbb{P}(\cdot)$ defined over $F$ is called 
a probability measure if it satisfies the following conditions:
\begin{enumerate}
    \item[$K1$] (Axiom of non-negativity) $\mathbb{P}(E) \geq 0 \quad \forall E \in F$. 
    \item[$K2$] (Axiom of unit norm) $\mathbb{P}(\Omega) = 1$.
    \item[$K3$] (Axiom of countable additivity) Any countable 
    sequence of disjoint sets $E_1, E_2, \ldots$ in $F$ satisfies 
    \begin{equation}
        \mathbb{P} \left( \bigcup^{\infty}_{i=1} E_i \right) = \sum^\infty_{i=1} \mathbb{P}[E_i].
    \end{equation}
\end{enumerate}
\end{axiom}

\begin{theorem}
    Let $F$ be a $\sigma$-field of events of sample space $\Omega$.
    \begin{enumerate}
        \item $\mathbb{P}(\varnothing) = 0$, where $\varnothing$ is empty set.
        \item $\mathbb{P}(E) \leq 1$, for any event $E \in F$. 
    \end{enumerate} 
\end{theorem}
\begin{proof}
    In the notion of Kolmogorov's axioms, the probability is interpret as a set function.
    \begin{enumerate}
        \item Let $E_1, E_2, \ldots$ be a sequence of disjoint events in $F$. Consider 
        $E_i = \varnothing$ for all $i$. Then 
        \[
            \bigcup^\infty_{i=1} E_i = \varnothing,
        \]
        and since $E_i \cap E_j = \varnothing \cap \varnothing =\varnothing, \forall i \neq j$. Hence 
        $E_i$'s are mutually exclusive.

        Now, with the axiom of countable additivity, we have
        \[
            \mathbb{P} \left( \bigcup^{\infty}_{i=1} E_i \right) = \sum^\infty_{i=1} \mathbb{P}[E_i].
        \]
        Or,
        \[
            \mathbb{P}(\varnothing) = \mathbb{P}(\varnothing) + \mathbb{P}(\varnothing) + \cdots.
        \]
        But this can only happens if either $\mathbb{P}(\varnothing) = 0$, $\mathbb{P}(\varnothing) = -\infty$, or 
        $\mathbb{P}(\varnothing) = +\infty$. Since $\mathbb{P}$ is a real valued function, so 
        $\mathbb{P}(\varnothing) = -\infty$ and $\mathbb{P}(\varnothing) = +\infty$ are absurd. 
        The only possible case is $\mathbb{P}(\varnothing) = 0$. This completes the proof.

        \item As $E \subset \Omega$ for each $E \in F$. We have 
        \[
            \mathbb{P}(E) \leq \mathbb{P}(\Omega). \label{eq:p1.0} \tag{{\color{red} $\heartsuit $}}
        \]
        From the axiom of unit norm, we knew $\mathbb{P}(\Omega) = 1$. Hence, \eqref{eq:p1.0} leads to
        $\mathbb{P}(E) \leq 1$ for any event $E$. Now we are done.

        \vspace{1.45cm}
        \textbf{Alternative proof}
        
        Notice that $E \cup E^c = \Omega$, and $E \cap E^c = \varnothing$.
        This implies $E \subset E^c$. By the law of finite additivity in probability measure $\mathbb{P}$,
        \begin{align*}
            \mathbb{P}(E \cup E^c) &= \mathbb{P}(E) + \mathbb{P}(E^c) - \mathbb{P}(E \cap E^c)\\
            \implies \mathbb{P}(\Omega) &= \mathbb{P}(E) + \mathbb{P}(E^c).
        \end{align*}
        Thus, we have $0 \leq \mathbb{P}(E^c) = 1 - \mathbb{P}(E)$ by Kolmogorov's axiom I. In result, 
        $\mathbb{P}(E) \leq 1$.
    \end{enumerate}
\end{proof}

\begin{theorem}[Total Probability]
    Let $(\Omega, F, \mathbb{P})$ be the probability space. Suppose 
    $\{H_n\}$ is a sequence of mutually exclusive and exhaustive events 
    such that $\mathbb{P}(H_n) > 0, H_n \in F$ for all $n \in \mathbb{N}$. 
    Then the probability of any event $A$ in $F$ is given by 
    \begin{equation}
        \mathbb{P}(A) = \sum^\infty_{n=1} \mathbb{P}(H_n)\, \mathbb{P}(A \mid H_n).
    \end{equation}
\end{theorem}